\documentclass[12pt]{article}
\usepackage[margin=1in]{geometry}
\usepackage{enumitem}
\usepackage{titlesec}
\usepackage{parskip}
\usepackage{amsmath}
\usepackage{hyperref}
\usepackage{fontenc}

\title{\textbf{Nalepa 2010 Claude Guide}\\
\large Nalepa, Monika. \textit{Skeletons in the Closet: Transitional Justice in Post-Communist Europe}.\\
New York: Cambridge University Press, 2010.}
\author{}
\date{}

\begin{document}

\maketitle
\tableofcontents
\newpage

%--------------------------------------------------
\section{Central Claims}
%--------------------------------------------------

\begin{itemize}[leftmargin=*, itemsep=4pt]
    \item The book's puzzle: why do some post-communist democracies in Central and Eastern Europe pursue extensive \textbf{lustration}---the vetting and public exposure of politicians and citizens who secretly collaborated with the communist-era secret police---while others pursue little or none, and why does the \emph{timing} and \emph{intensity} of lustration vary so dramatically across cases that otherwise look similar (all newly democratic, all emerging from communist one-party rule, all facing a common legacy of secret police files)?
    \item Nalepa's answer rejects both a purely normative account (``some societies simply value justice/truth more'') and a purely structural account (regime-type-at-exit determines outcomes) in favor of a \textbf{strategic, game-theoretic logic}: lustration outcomes are the equilibrium result of politicians' incentives given what they know about their own and others' collaboration histories, not simply a society-wide referendum on dealing with the past.
    \item The titular metaphor---\textbf{skeletons in the closet}---captures the core mechanism: some politicians on \emph{both} sides of the old regime/opposition divide have hidden pasts as secret collaborators. These politicians (whether nominally ex-communist or nominally former dissidents) have private incentives to block, weaken, delay, or narrow lustration to avoid having their own skeleton exposed.
    \item Conversely, politicians \emph{without} skeletons have incentives to support lustration both for principled/programmatic reasons (retributive or truth-seeking justice) and for strategic reasons: exposing a rival's collaborationist past is a way to discredit and eliminate electoral competition. Lustration is therefore never purely backward-looking justice; it is simultaneously a \textbf{forward-looking political weapon}.
    \item A key insight is that support for or opposition to lustration cannot be read directly off party labels (ex-communist vs.\ opposition) because collaboration was not confined to the communist party--secret police apparatus alone; opposition and dissident circles were themselves infiltrated by informers. This scrambles the naive expectation that ex-communists always oppose lustration and former dissidents always support it.
    \item The variation Nalepa seeks to explain is not simply ``lustration happened vs.\ did not,'' but finer-grained: \emph{when} lustration laws were passed relative to the founding transitional election, how broad their scope was (which offices/professions were covered), how the process was administered (courts vs.\ special commissions), and whether secret police archives were preserved, sealed, or destroyed at the moment of transition.
    \item The central theoretical move is to treat the \textbf{transition pact itself}---the founding-moment bargain between outgoing communist elites and incoming opposition negotiators---as the pivotal strategic moment that locks in downstream possibilities for lustration, because it is at the pact that decisions about the fate of the secret police archives (destroy, hide, preserve intact) are made.
    \item Because archive preservation and lustration timing are chosen strategically by actors who already know (probabilistically) their own skeleton status, the empirical correlation between ``who negotiated the pact'' and ``how extensive lustration later became'' is not incidental---it is the observable footprint of the underlying asymmetric-information bargaining game.
\end{itemize}

%--------------------------------------------------
\section{Core Theoretical Model}
%--------------------------------------------------

\begin{itemize}[leftmargin=*, itemsep=4pt]
    \item Nalepa formalizes the founding moment of democratic transition as a \textbf{bargaining game under asymmetric information} between two blocs of actors: outgoing regime elites (the communist nomenklatura) and incoming opposition leaders who will inherit government. Both blocs are internally heterogeneous with respect to a hidden type: whether a given individual politician was or was not a secret police collaborator/informer.
    \item \textbf{Private information}: each politician knows his or her own collaboration history with certainty (or near-certainty) but does not know with certainty the collaboration histories of other politicians, including nominal allies. This is the asymmetric-information core of the model---collaboration status is a privately known ``type.''
    \item \textbf{The strategic role of the secret police archives}: the files are the only source of hard, verifiable evidence about who collaborated. Whether they are destroyed, hidden, or preserved at the transition moment is therefore not a neutral administrative detail but a decision with first-order consequences for future political competition---it determines whether skeletons can ever be verifiably exposed.
    \item \textbf{Signaling/screening logic around lustration}: because a politician's willingness to support strict, verifiable lustration is correlated with (though not dispositive proof of) that politician's clean type, support for or opposition to lustration functions as an (imperfect, costly) signal of collaboration status. Politicians with skeletons prefer either no lustration or a form of lustration weak/manipulable enough that they can survive it (e.g., self-declaration systems without independent verification, narrow scope, generous statutes of limitations).
    \item \textbf{Lustration as a screening device wielded strategically}, not (only) a retributive institution: incoming democratic politicians without skeletons have an incentive to design lustration rules that will differentially expose rivals--- turning transitional justice into a tool of intra- and inter-party competition rather than a straightforwardly moral reckoning with the past.
    \item \textbf{Why pacts vary in their treatment of the archives}: at the pact-negotiation stage, the outgoing regime bloc contains a mix of collaborators and (relatively) clean elites, as does the incoming opposition bloc. The equilibrium outcome for the fate of the archives (and hence for later lustration) depends on the relative bargaining power and the distribution of skeleton-types within each bloc at that founding moment---not simply on which side ``won'' the transition.
    \item \textbf{Commitment problem}: outgoing elites cannot credibly commit to a lenient future if they hand over/preserve incriminating files, and incoming opposition politicians with skeletons cannot credibly commit to supporting harsh lustration once in power. This generates a structural parallel to bargaining problems over commitment in other transition and power-sharing contexts (see Connections, below, on Svolik 2012).
    \item \textbf{Equilibrium predictions}: the model generates testable comparative statics---e.g., pacts negotiated by blocs with a higher proportion of skeleton-bearing elites on both sides are more likely to result in destroyed/inaccessible archives and weak, delayed, or narrow lustration; pacts negotiated by relatively ``clean'' opposition blocs facing a compromised or fractured regime bloc are more likely to produce early, broad, and verifiable lustration.
\end{itemize}

%--------------------------------------------------
\section{Core Empirical Strategy}
%--------------------------------------------------

\begin{itemize}[leftmargin=*, itemsep=4pt]
    \item \textbf{Comparative case studies across Central and Eastern Europe}, principally Poland, the Czech Republic (and former Czechoslovakia), and Hungary, supplemented by comparisons to other post-communist cases (e.g., East Germany/the former GDR via the Stasi archive experience, and briefer treatment of other regional cases) to trace variation in lustration scope, timing, and intensity.
    \item \textbf{Poland}: treated as the paradigm case of a negotiated ``round table'' pact (1989) between the communist regime and the Solidarity-led opposition. Lustration in Poland was notably delayed, contested, and repeatedly watered down/relitigated over the 1990s and 2000s (e.g., the Kwa\'sniewski-era controversies, the eventual 1997 lustration law and its troubled implementation, and periodic later re-openings such as the 2006/2007 broadening under the Kaczy\'nski government). Nalepa reads this halting, repeatedly-revised trajectory as consistent with a founding pact negotiated by blocs on both sides containing significant numbers of politicians with skeletons.
    \item \textbf{Czech Republic/Czechoslovakia}: contrasted as a case of comparatively early, broad, and strict lustration (the 1991 lustration law), read as reflecting a transition in which the incoming opposition (Civic Forum) had comparatively fewer skeleton-bearing elites and/or greater relative bargaining leverage over a discredited, non-negotiating regime (the ``Velvet Revolution'' was less of a negotiated pact and more of a regime collapse), making broad lustration a more attainable and less personally risky equilibrium for incoming politicians.
    \item \textbf{Hungary}: treated as an intermediate case---a negotiated transition (the Hungarian Round Table) but with distinctive patterns in archive access and lustration scope/timing that Nalepa uses to further test the pact-composition logic, including protracted controversies over ``Duna-gate''-type revelations and repeated legislative revisitations of lustration rules.
    \item \textbf{Data and evidence}: the empirical chapters combine (a) qualitative process-tracing of the transition negotiations and subsequent legislative histories of lustration laws in each country, (b) an accounting of which negotiators/parties are later revealed (through archives, memoirs, journalism, or subsequent lustration proceedings) to have had collaboration histories, and (c) in Nalepa's broader research program, original survey and archival data (including list experiments and other indirect measurement techniques designed to elicit sensitive information about collaboration and attitudes toward lustration that respondents would not reveal to direct questioning).
    \item \textbf{Identification strategy}: the core empirical test is whether the \emph{composition of the pact-negotiating elites' skeleton status}---not merely regime type, level of communist repression, or economic development---predicts the subsequent scope/timing/intensity of lustration and the fate of the secret police archives. Cross-case comparison is used because within any single case the key explanatory variable (private collaboration status of negotiators at the time of the pact) is not directly observable to the researcher either, mirroring the informational problem inside the model itself.
    \item \textbf{Robustness checks/extensions}: Nalepa examines within-case variation over time (e.g., repeated Polish lustration law revisions) as a way of tracing how newly revealed information about individual politicians' skeleton status shifts the political coalition supporting or opposing further lustration, providing dynamic, over-time evidence beyond the static cross-national comparison.
\end{itemize}

%--------------------------------------------------
\section{Core Works Cited}
%--------------------------------------------------

\begin{itemize}[leftmargin=*, itemsep=4pt]
    \item \textbf{Guillermo O'Donnell and Philippe Schmitter}, \textit{Transitions from Authoritarian Rule: Tentative Conclusions about Uncertain Democracies} (1986)---the foundational framework for ``pacted transitions'' that Nalepa directly engages and formalizes; her bargaining-game model of the transition moment is best read as a game-theoretic microfoundation for the O'Donnell/Schmitter account of elite pact-making.
    \item \textbf{Adam Przeworski}, work on democratic transitions and elite pacts (e.g., \textit{Democracy and the Market}, 1991)---shared rational-choice institutionalist tradition; Przeworski's treatment of transitions as strategic games among elites with uncertain payoffs under the new regime is a direct intellectual ancestor of Nalepa's model.
    \item \textbf{Kathryn Sikkink}, work on the ``justice cascade'' and human rights prosecutions (e.g., \textit{The Justice Cascade}, 2011)---broader transitional-justice literature situating truth commissions and prosecutions comparatively; provides context for how post-communist lustration fits into (and differs from) the global spread of transitional justice mechanisms.
    \item \textbf{Ruti Teitel}, \textit{Transitional Justice} (2000)---a foundational typology of transitional justice mechanisms (retributive, reparative, reconstructive); Nalepa's lustration-as-strategic-weapon account complicates Teitel's more normatively oriented framework by emphasizing the self-interested strategic uses of transitional justice institutions.
    \item \textbf{Timothy Garton Ash}, journalistic/historical accounts of Central European lustration debates and secret police archives (e.g., \textit{The File: A Personal History}, 1997, and essays on Central European lustration)---provides rich descriptive/historical texture on the Czech, Polish, and German lustration controversies that Nalepa's formal model seeks to explain more systematically.
    \item \textbf{Alexandra Barahona De Brito, Carmen Gonz\'alez-Enr\'iquez, and Paloma Aguilar} (eds.), \textit{The Politics of Memory: Transitional Justice in Democratizing Societies} (2001)---comparative volume situating post-communist lustration alongside Latin American and Southern European transitional justice experiences.
    \item \textbf{Formal/game-theoretic transitions literature} more broadly (e.g., work on signaling games and screening in political economy)---the general toolkit of asymmetric-information bargaining models that Nalepa adapts to the specific institutional setting of secret police files and lustration.
\end{itemize}

%--------------------------------------------------
\section{Methodological Contributions and Limitations}
%--------------------------------------------------

\begin{itemize}[leftmargin=*, itemsep=4pt]
    \item \textbf{Signature methodological contribution}: the tight integration of a formal game-theoretic model (asymmetric information, signaling/screening over hidden collaboration types) with detailed comparative-historical case evidence (Poland, Czech Republic, Hungary) and, in Nalepa's broader research program, original micro-level data (archival and survey evidence, including indirect elicitation techniques such as list experiments designed to measure sensitive attitudes toward lustration and collaboration without direct self-incrimination).
    \item \textbf{Strength}: the formal model generates precise, falsifiable comparative statics (e.g., pact composition predicts archive fate and lustration scope) rather than relying on looser narrative claims about ``political will'' or ``civic culture,'' giving the comparative case studies clear theoretical discipline.
    \item \textbf{Core empirical limitation}: the model's key causal variable---individual politicians' private knowledge of their own and others' collaboration status at the moment of the pact---is inherently difficult to verify empirically. Researchers (like the politicians in the model) typically cannot directly observe who was secretly informing at the time negotiations occurred; evidence tends to surface only later, if at all, through subsequent archive openings, making rigorous testing of the model's micro-foundations harder than testing its aggregate predictions.
    \item \textbf{Selection and archive-survival concerns}: because the very object of study---secret police files---was itself destroyed, altered, or selectively preserved by strategic actors, the empirical record available to test the theory is not a random sample of the truth; missing or doctored archives can bias inferences about who actually had ``skeletons.''
    \item \textbf{Generalizability limits}: the mechanism is tightly tailored to a specific institutional feature of post-communist Europe---the existence of extensive, individualized secret police surveillance archives naming citizen-informers. Extending the logic to other transitional justice settings (e.g., military dictatorships in Latin America, apartheid South Africa) requires care, since those settings often lack comparably extensive, individually incriminating written archives, changing the informational structure of the underlying game (verification of ``skeletons'' may be impossible rather than merely difficult).
    \item \textbf{Scope conditions}: the theory presumes a pacted (negotiated) rather than a collapsed/imposed transition is at least a live possibility, since the bargaining logic hinges on both outgoing and incoming elites having some leverage at the founding moment; cases of abrupt regime collapse without negotiation (arguably closer to the Czechoslovak case) sit somewhat differently within the framework than fully negotiated round-table pacts (Poland, Hungary).
    \item \textbf{Contribution to the transitional justice literature}: reframes lustration away from a purely normative or moralized debate (justice vs.\ reconciliation) and toward a positive political-economy account of institutional design under uncertainty, showing that seemingly principled positions on transitional justice can be strategically motivated equilibrium behavior.
\end{itemize}

%--------------------------------------------------
\section{Connections to the Broader Literature}
%--------------------------------------------------

\begin{itemize}[leftmargin=*, itemsep=4pt]
    \item \textbf{Przeworski et al.\ 2000 (\textit{Democracy and Development})}: shares the rational-choice, game-theoretic approach to democratic transitions and elite strategic behavior around regime change; both treat transition outcomes as equilibria of strategic interaction among self-interested elites rather than as simple products of structural or cultural preconditions.
    \item \textbf{Huntington 1968}: Huntington's concern with the destabilizing effects of rapid political change and the challenge of building durable post-transition institutions bears directly on Nalepa's account---the terms struck at the pact (including the fate of the secret police archives) shape the institutional trajectory and legitimacy of the new democratic regime for years afterward, echoing Huntington's emphasis on institutionalization as the key variable determining whether rapid change produces order or instability.
    \item \textbf{Svolik 2012 (\textit{The Politics of Authoritarian Rule})}: Svolik's analysis of the commitment problems between dictators and their ruling coalitions/inner circle is close kin to Nalepa's model of the commitment problem between outgoing authoritarian elites and incoming democratic opposition negotiators---both are fundamentally about actors who cannot credibly bind their future selves (or successors) regarding the treatment of information/power once circumstances change.
    \item \textbf{Moore 1966 and Boix 2003}: both offer class/elite-centered structural theories of regime transition (Moore's ``routes to modernity'' via class coalitions; Boix's inequality/redistribution-based model of democratization). These are useful contrasts to Nalepa's more institution- and information-centered strategic account, which focuses specifically on the moment and terms of transition (and the informational structure of the pact) rather than on long-run structural or distributional preconditions for democratization.
    \item \textbf{Weber 1968}: Weber's account of legal-rational authority and the routinization/institutionalization of legitimate rule is relevant to how a new democratic regime's chosen stance on lustration (embracing, evading, or manipulating transitional justice) shapes its relationship to---and legitimation vis-\`a-vis---the authoritarian past, and to whether the new regime's authority is grounded in a clean break or a compromised continuity with the old elite.
    \item \textbf{Ziblatt 2017} (where referenced): Ziblatt's account of conservative elites' organizational capacity and commitment across a regime transition offers a parallel elite-centered logic of how founding-moment elite behavior shapes long-run regime trajectories, though Ziblatt's focus is on party-building capacity rather than on hidden information/skeletons as the driving strategic variable.
    \item \textbf{Kalyvas 2006 (\textit{The Logic of Violence in Civil War})}: Kalyvas's analysis of information, denunciation, and collaboration under conditions of political violence and contested control is structurally analogous to Nalepa's problem---both concern actors possessing hidden information about who collaborated with a controlling power, the strategic incentives to reveal or conceal that information, and how the local/institutional structure for verifying collaboration (zones of control for Kalyvas; secret police archives for Nalepa) shapes political outcomes. Kalyvas's collaboration/denunciation logic under coercive control is arguably the closest structural parallel among the readings to Nalepa's core informational mechanism.
    \item \textbf{Cox 1987 (\textit{The Efficient Secret})}: while focused on a different institutional question (the emergence of disciplined parliamentary parties in Britain), Cox's attention to how institutional rules shape the strategic incentives of political actors resonates with Nalepa's broader claim that seemingly moral/normative political stances (like support for lustration) can be reinterpreted as strategically rational responses to institutional incentive structures.
    \item \textbf{Broader transitional justice literature (Sikkink, Teitel, Garton Ash)}: situates the post-communist lustration cases within the global rise of transitional justice mechanisms, while Nalepa's game-theoretic reframing pushes back against purely norm-diffusion or moral-progress narratives (e.g., Sikkink's ``justice cascade'') by showing how strategic self-interest around hidden pasts can drive---or block---the adoption of transitional justice institutions independent of normative consensus.
\end{itemize}

\end{document}
